\documentclass[11pt]{article}
\usepackage[margin=1in]{geometry}
\usepackage{amsmath,amssymb,amsthm,mathtools}
\usepackage{graphicx,xcolor,enumitem,microtype,hyperref}
\hypersetup{colorlinks=true,linkcolor=blue!55!black,citecolor=blue!55!black,urlcolor=blue!55!black}
\newtheorem{theorem}{Theorem}
\newtheorem{lemma}{Lemma}
\newcommand{\A}{\mathcal A}
\newcommand{\T}{\mathbb T}
\newcommand{\wind}{\operatorname{wind}}
\title{\vspace{-2em}A Cubic with No Critical Element over $C(S^1)$\\
\large Counterexample to two C*-algebraic mean value conjectures}
\author{Counterexample packet for arXiv:2206.08154}
\date{26 August 2026}

\begin{document}
\maketitle
\begin{center}
\fcolorbox{red!65!black}{red!3}{\parbox{0.91\textwidth}{
\textbf{Status: candidate full counterexample, likely valid.}
There is a degree-three polynomial over the commutative C*-algebra
$C(S^1)$ whose derivative has no zero in the algebra.  At an explicitly
chosen noncritical point, this contradicts the existential conclusion common
to the C*-algebraic Smale and Dubinin--Sugawa conjectures, including their
strong forms.
}}
\end{center}

\section{Source and exact conjectures}

The source is K.~Mahesh Krishna,
\emph{C*-algebraic Smale Mean Value Conjecture and Dubinin--Sugawa Dual Mean
Value Conjecture}, arXiv:2206.08154 (2022) \cite{Krishna}.
For a commutative C*-algebra $\A$ and
\(P(x)=\prod_{j=1}^n(x-a_j)\), it defines a critical element to be
$w\in\A$ with $P'(w)=0$.

Conjecture~2.1 says that if $z$ is not critical, there is a critical
$w\in\A$ satisfying the proposed Smale upper estimate.  Conjecture~3.1
says that, under the same hypotheses, there is a critical $w\in\A$ satisfying
the proposed Dubinin--Sugawa lower estimate.  Conjectures~2.2 and 3.2 make
stronger C*-order assertions but retain the same existence claim.  The formal
statements are on source pages 4--5 and 7; both complete conjectures also
appear in the abstract on source page 1.

\begin{figure}[p]
\centering
\includegraphics[width=0.98\textwidth,height=0.65\textheight,keepaspectratio]{figures/open_problem_crop.png}
\caption{The source abstract on PDF page 1, containing the two conjectures in
full and noting that the paper proves only degree two.}
\end{figure}
\clearpage

\section{Proof intuition}

For a cubic, solving $P'(w)=0$ requires choosing a square root of its
discriminant.  Over the scalar field this is automatic.  Over $C(S^1)$, an
element is a continuous function, so a root of the derivative must amount to
a continuous choice of one of the two scalar critical points in every fiber.

We arrange for the derivative discriminant to trace a nonvanishing loop of
odd winding number.  It has two distinct square roots in every individual
fiber but no continuous square root around the circle: traversing the circle
exchanges the two roots.  Hence the polynomial has no critical element at
all.  Choosing $z=0$ with $P'(0)$ invertible meets the conjectures' hypothesis,
while their required $w$ cannot exist; the numerical inequalities never enter.

\section{The construction and proof}

Write $\T=\{\zeta\in\mathbb C:|\zeta|=1\}$ and let
$u\in C(\T)$ be the coordinate function $u(\zeta)=\zeta$.  Set
\[
 \rho=e^{i\pi/3}=\frac{1+i\sqrt3}{2},\qquad
 \varepsilon=\frac14,\qquad c=\rho+\varepsilon u.
\]

\begin{lemma}[Odd winding has no continuous square root]\label{lem:winding}
If $q\in C(\T)$ is nowhere zero and $\wind(q)$ is odd, then there is no
$h\in C(\T)$ with $h^2=q$.
\end{lemma}

\begin{proof}
Such an $h$ would also be nowhere zero.  Winding number is additive under
pointwise multiplication, so
$\wind(q)=\wind(h^2)=2\wind(h)$, an even integer.
\end{proof}

\begin{theorem}\label{thm:counterexample}
In the unital commutative C*-algebra $\A=C(\T)$, the cubic
\[
                    P(x)=x(x-1)(x-c)
\]
has no critical element.  Nevertheless $z=0$ is not critical.  Consequently
Conjectures~2.1, 2.2, 3.1, and 3.2 of \cite{Krishna} are false in degree
three.
\end{theorem}

\begin{proof}
Expanding in the commutative algebra gives
\[
 P(x)=x^3-(1+c)x^2+cx,\qquad
 P'(x)=3x^2-2(1+c)x+c.
\]
At $z=0$ we have $P'(0)=c$.  Pointwise,
\[
              |c(\zeta)|=|\rho+\varepsilon\zeta|
                 \geq 1-\varepsilon=\frac34,
\]
so $c$ is invertible in $C(\T)$; in particular, $z$ is not critical.

Suppose that $w\in C(\T)$ were a critical element.  From $P'(w)=0$ we get
\[
 9w^2=6(1+c)w-3c.
\]
Thus, with $h=3w-(1+c)$,
\begin{align*}
 h^2
 &=9w^2-6(1+c)w+(1+c)^2 \\
 &=(1+c)^2-3c=c^2-c+1.
\end{align*}
Because $\rho^2-\rho+1=0$ and $2\rho-1=i\sqrt3$,
\begin{equation}\label{eq:qfactor}
 q:=c^2-c+1
   =\varepsilon u\bigl(i\sqrt3+\varepsilon u\bigr).
\end{equation}
The second factor never vanishes since
$|i\sqrt3+\varepsilon\zeta|\geq\sqrt3-\varepsilon>0$.
Its image is a circle centered at $i\sqrt3$ of radius $\varepsilon$, which
does not surround the origin, so its winding number is zero.  The factor
$\varepsilon u$ has winding number one.  Therefore
\[
                         \wind(q)=1+0=1.
\]
Equation \eqref{eq:qfactor} and Lemma~\ref{lem:winding} contradict $h^2=q$.
Hence $P'$ has no zero in $C(\T)$.

The polynomial has degree three and the required factor form with roots
$0,1,c\in\A$.  Each of the four cited conjectures applies at a noncritical
$z$ and concludes that some critical $w\in\A$ exists.  Our admissible
$z=0$ has no such $w$, so all four conclusions fail.
\end{proof}

\paragraph{Fiberwise interpretation.}
The obstruction is genuinely global, not scalar.  At every
$\zeta\in\T$, the scalar derivative has the two distinct roots
\[
 w_\pm(\zeta)=\frac{1+c(\zeta)\pm\sqrt{q(\zeta)}}3,
\]
because $q(\zeta)\neq0$.  Odd monodromy prevents either choice from defining
a continuous function on the whole circle.

\section{Verification, novelty, and scope}

\paragraph{Verification.}
The argument is exact and uses only polynomial algebra and the elementary
winding-number obstruction.  The included checker samples 16,384 points,
verifies \eqref{eq:qfactor}, the lower bounds for $|c|$ and $|q|$, winding
number one, and the scalar derivative roots.  It is not used as proof.

\paragraph{Bounded novelty search.}
The four run indexes were searched for arXiv:2206.08154, the exact conjecture
titles, $C(S^1)$, critical points, continuous square roots, monodromy, and
winding.  Exact-title, exact-phrase, author/id, counterexample, degree-three,
and close-variant arXiv/web searches through 26 August 2026 found the source
paper but no later paper claiming this construction or an exact resolution.
Novelty is plausible, not certified.

\paragraph{Scope and review.}
This is a full counterexample to the unrestricted commutative C*-algebra
statements in Conjectures~2.1 and 3.1 and, a fortiori, to their proposed
strong forms.  It does not address the classical scalar Smale or
Dubinin--Sugawa conjectures.  It also leaves open versions over algebras whose
spectrum has no square-root monodromy obstruction.  Human review should check
the derivative identity, the source's critical-point definition, and the
winding calculation.

\begingroup
\sloppy
\raggedright
\begin{thebibliography}{9}
\bibitem{Krishna}
K.~Mahesh Krishna,
\emph{C*-algebraic Smale Mean Value Conjecture and Dubinin--Sugawa Dual Mean
Value Conjecture},
\href{https://arxiv.org/abs/2206.08154}{arXiv:2206.08154}, 2022.
\end{thebibliography}
\endgroup
\end{document}
